\section{Introduction}
\label{sec:intro}

Large language model (LLM) serving has moved from single-GPU engines to
\emph{disaggregated} deployments in which prefill and decode run on different workers, KV cache
travels between them, and each incoming request must be routed to \emph{some} worker
\cite{distserve,splitwise,mooncake,llumnix}. Routing is no longer a detail: the worker whose
compute queue is emptiest is frequently \emph{not} the worker whose link can deliver the
request's context (or its KV cache) fastest, and the two criteria can point in opposite
directions. A growing line of work therefore makes the router \emph{network-aware}: it scores
candidates using both a compute-backlog term and a link term
\cite{cassini,vclos,helix,netkv,aegis,gensched}.

The problem is that this literature reports wildly different value for the same idea.
Reported gains range from negligible to more than $70\%$, and they are usually quoted as a
single number under a single configuration. Worse, at least one recently published system
reports a property that ought to be alarming for anyone deploying this class of scheduler:
its gains ``hold under workload shifts and \emph{telemetry delays up to 60 seconds}''
\cite{aegis}. If stale network telemetry does not matter, then a large fraction of the
engineering effort around network-aware routing---and its sensitivity to monitoring
latency---is misdirected. Yet the same paper, and the closest concurrent work
\cite{netkv}, sweep only \emph{one} observable quantity.

This paper asks the prior question: \textbf{when does network awareness actually pay in LLM
inference serving, and when does trusting it hurt?} We treat the scheduler's benefit as an
outcome to be \emph{characterized} rather than a result to be demonstrated: we measure where it
appears, where it vanishes, where it reverses sign, and how each factor interacts with the
others. Three design decisions shape the study.

\paragraph{Ten independent workload realizations, not one.}
A single simulation of a congested queueing system is not an experiment. In preliminary
single-seed screens we observed that changing the routing decisions of the first $\sim2.5\%$ of
requests alone moved the 95th-percentile sojourn time by up to $17\%$ in a 4000-request
run---i.e.~single-seed comparisons of this system are dominated by trajectory noise. We therefore
treat the \emph{seed}, i.e.~one independent realization of the arrival process, as the unit of
observation: every comparison in this paper aggregates ten seeds, uses per-request paired tests
within a seed, and quotes intervals over seeds. Comparisons that a single-seed screen would have
reported as ``network awareness becomes harmful under stale state'' turn out to be artifacts of
one extreme seed (we give the concrete example in \S\ref{sec:discount}).

\paragraph{One mechanism variable, measured directly.}
Instead of reporting that network awareness ``helps'', we identify a configuration variable that
strongly modulates---under the controlled conditions of our experiments---whether the link term
carries \emph{any} information beyond the compute term: the
rank correlation between per-worker compute capacity and per-worker bandwidth (that is, whether
the fastest-computing workers are also those with the widest links). When the
fastest-computing worker also has the widest link, the two criteria select the same worker and
the link term is redundant; when they disagree, the link term is the only thing standing between
the user and a $68\%$ worse tail. This distinction---which we call \emph{alignment}---is, in our
reading of the literature, absent as an explanatory variable; within our fixed compute sets,
bandwidth sets, workloads, and routing rules it accounts for a large part
of the variance in reported gains, while absolute capacity ratios, topology, payload size, and
queueing discipline remain co-determinants we do not fully separate.

\paragraph{Staleness is not one thing but two, with different expiry times.}
The scheduler observes two physically different quantities: the compute queue at each worker
(changes on sub-second to second timescales) and the state of each link (bandwidth is a physical
property, congestion changes on tens of seconds). Each observation carries an \emph{age}---how
many seconds old the message is when the scheduler acts on it---and we sweep them jointly and find their value
expires at rates that differ by roughly an order of magnitude. This single fact explains why
the published ``60\,s does not matter'' claim is both true and misleading: it was measured on the
slow-expiring channel only. We further separate \emph{latency} (seeing an old but true value)
from \emph{averaging} (seeing a time-averaged value, where spikes are smoothed away), and show
the two distort decisions differently.

\paragraph{Contributions.}
\begin{itemize}
  \item \textbf{An explanatory variable.} We introduce compute--bandwidth rank alignment as an
        important explanatory factor---among several, under our experimental controls---for
        whether network information adds anything, and validate it across three
        nodes $\times$ ten seeds (effect separates completely: Cliff's $\delta=1.00$,
        $p=1.8\times10^{-4}$), across 4/8/16 workers, and against a direct probe that rules out
        the obvious alternative explanation.
  \item \textbf{A two-channel model of staleness with measured expiry times,} plus a 2-D
        (queue age $\times$ link age) sweep, plus a separate characterization of averaging
        windows as a distinct distortion from latency.
  \item \textbf{A conditional discounting rule} whose parameters are fitted from measurement
        rather than tuned, together with the \emph{falsification} of the intuitive alternative
        and a three-source sensitivity ablation. We also show the rule is insurance, not a
        free win: it is a net loss outside one regime, and we say which.
  \item \textbf{Measurement of the observation feedback loop on real containers,} which the
        simulator abstracts away: fetching state costs \SIrange{56}{876}{\milli\second} against
        12 µs for an in-cluster round trip, and it grows with load. The
        tail-latency impact is threshold-like rather than proportional.
  \item \textbf{A reproducibility protocol and an explicit negative-results ledger.} Four
        validation gates, append-only artifacts, a manifest per run, a script that re-derives
        every number in the paper from the raw outputs, and a section that lists what did not
        work---including a design of our own that we falsify with our own data.
\end{itemize}

We are explicit about what this paper is not. It does not propose a scheduler that beats the
state of the art; our conditional discounting rule is, in several regimes, indistinguishable
from simply always trusting the network state (\S\ref{sec:discount}). It does not evaluate a
multi-node GPU cluster: our real-cluster evidence comes from a single-node \texttt{k3s}
deployment with application-level link shaping, which we use for direction validation and for
quantities the simulator omits. And it does not fully explain when discounting pays: we tested
two candidate mechanisms (arrival burstiness and rate persistence) and \emph{both were refuted
by our own data}; we report that as an open question rather than fitting a story to it.
